\section{Manus: la ingeniería de contexto es el núcleo de un Agent en producción}

Kimi, ChatGPT Agent y Qwen muestran, respectivamente, RL aplicado al modelo, al producto y al código; esta sección examina la capa más subestimada de un Agent en producción: la ingeniería de contexto. Manus optó por construir su Agent sobre la capacidad in-context de los modelos de frontera, en lugar de entrenar desde cero un modelo de extremo a extremo. Subraya que la ingeniería de contexto es una ciencia experimental y que lo decisivo en un Agent en producción no es solo el prompt, sino el KV cache, la gestión de herramientas, la memoria en el sistema de archivos, la reiteración del todo, la conservación de los errores y evitar la trampa de los pocos ejemplos.

\begin{figure}[H]
\centering
\includegraphics[width=0.96\textwidth]{manus-context-engineering.png}
\caption{Manus Context Engineering: KV cache, enmascaramiento de herramientas, sistema de archivos, todo, conservación de errores y diversidad. Diagrama conceptual propio, elaborado a partir de la publicación oficial del blog de Manus y de la conversación de 01:59:04--02:06:06.}
\end{figure}

\begin{knowledgebox}{Lectura de la figura: la ingeniería de contexto es ingeniería en tiempo de ejecución}
El entrenamiento le da capacidades al modelo; la ingeniería de contexto permite que el Agent las use en producción de forma estable, económica y recuperable. Se ocupa de cada ronda de llamadas a herramientas, del crecimiento del contexto y de la recuperación ante fallos.
\end{knowledgebox}

\subsection{KV cache y enmascaramiento de herramientas}

La sección anterior planteó que la ingeniería de contexto es ingeniería en tiempo de ejecución; esta desglosa los dos principios más prácticos: el KV cache y el enmascaramiento de herramientas. Manus considera que la tasa de aciertos del KV-cache es uno de los indicadores más importantes de un Agent en producción, porque la proporción entre entrada y salida de un Agent está muy cargada hacia la entrada, y la caché afecta directamente la latencia y el costo. También propone «enmascarar, no retirar» herramientas: agregar y quitar herramientas de forma dinámica rompe la caché y confunde al modelo; es mejor conservar las definiciones de las herramientas y controlar cuáles pueden elegirse mediante una máquina de estados o un logits mask.

\begin{knowledgebox}{Términos clave: puntos centrales de la ingeniería de contexto}
\begin{tabularx}{\textwidth}{p{0.22\textwidth}p{0.34\textwidth}X}
\toprule
Técnica & Función & Práctica errónea \\
\midrule
KV cache & Reduce el costo de repetir la inferencia sobre prefijos largos & Cambiar el prompt de sistema o las definiciones de herramientas en cada ronda. \\
Tool masking & Controla las acciones disponibles en el momento & Eliminar herramientas de forma dinámica, lo que vuelve inconsistente el contexto. \\
Filesystem as context & Usa el sistema de archivos como memoria externa de largo plazo & Meter todo el contenido en la ventana de contexto. \\
Todo.md & Dirige la atención mediante la reiteración de los objetivos & No actualizar los objetivos en tareas largas. \\
Keep errors & Conserva la evidencia de los fallos para facilitar la recuperación & Limpiar los errores, lo que lleva a repetirlos. \\
\bottomrule
\end{tabularx}
\end{knowledgebox}

\begin{importantbox}{Lista de verificación de un Agent en producción}
Un Agent en producción debe responder al menos seis preguntas: ¿el contexto puede almacenarse en caché? ¿El espacio de herramientas es controlable? ¿Dónde reside la memoria de largo plazo? ¿Se conservan los fallos? ¿Los objetivos se reiteran en el contexto reciente? ¿Los ejemplos repetidos pueden encasillar al modelo en un patrón? Casi toda la experiencia de Manus gira en torno a estas preguntas.
\end{importantbox}

\begin{knowledgebox}{Checklist del sistema de un Agent en producción}
\begin{tabularx}{\textwidth}{p{0.22\textwidth}p{0.34\textwidth}X}
\toprule
Elemento & Pregunta & Consecuencia de omitirlo \\
\midrule
Caché & ¿El prefijo es estable y el contexto es append-only? & El costo y la latencia se disparan. \\
Herramientas & ¿Las definiciones de herramientas son estables y las acciones pueden restringirse? & Elección de la herramienta equivocada, invalidación de la caché, alucinación de acciones. \\
Memoria & ¿El estado de largo plazo se externaliza en archivos o en una base de datos? & Explosión del contexto o pérdida de información. \\
Errores & ¿Las trayectorias fallidas se conservan para que el modelo aprenda de ellas? & Se repite el mismo error. \\
Seguridad & ¿Las acciones de alto riesgo se confirman y se auditan? & Una operación errónea del Agent causa daños reales. \\
\bottomrule
\end{tabularx}
\end{knowledgebox}

\subsection{Resumen de la sección}

Manus nos recuerda que los problemas de un Agent en producción muchas veces no se deben a que el modelo carezca de capacidad, sino a que el contexto, las herramientas, la memoria y la recuperación ante fallos no están bien diseñados.
